\documentclass[10pt]{article}
\usepackage[a4paper,margin=0.65in]{geometry}
\usepackage{amsmath,amssymb}
\usepackage{array}
\usepackage{enumitem}
\usepackage{xcolor}
\usepackage[hidelinks]{hyperref}
\setlist[itemize]{leftmargin=1.2em,itemsep=1pt,topsep=2pt}
\setlist[enumerate]{leftmargin=1.4em,itemsep=1pt,topsep=2pt}
\setlength{\parindent}{0pt}
\setlength{\parskip}{3pt}
\renewcommand{\arraystretch}{1.15}
\newcommand{\code}[1]{\texttt{#1}}

\begin{document}
\sloppy

%---------------- Title block ----------------
\begin{center}
  \rule{\linewidth}{0.6pt}\\[4pt]
  {\Large\bfseries DA3408 -- Assignment 5}\\[2pt]
  {\large SwiftParcel ETA Service -- Report}\\[2pt]
  \rule{\linewidth}{0.6pt}\\[6pt]
  \begin{tabular}{c}
    \textbf{Kiran Kumar P}\\
    DA24B008
  \end{tabular}
\end{center}

\section{Machine}
AMD Ryzen 7 8840HS (8 cores / 16 threads), 14\,GiB RAM, Ubuntu (Linux 7.0), Python 3.10.22 (conda).
fastapi 0.162.0, uvicorn 0.54.0, pydantic 2.14.3, scikit-learn 1.7.2, mlflow 3.18.0, ray 2.60.0, prometheus\_client 0.26.0.
The load generator ran on the same laptop as the server. All numbers come from runs on 5 Nov 2026 (raw outputs in
\code{partA/*.json}, \code{partB/*.txt}, \code{partC/verify\_output.txt}).

%================================================
\section{Part A: REST Service}

A1 and A2: \code{partA/eta\_app.py} and \code{partA/test\_eta\_app.py} (\textbf{15 tests pass}; breaking the
cross-field rule and the header echo in a copy of the app makes 4 fail).

\subsection{A3. Load Test (v1, batch 1, 4\,s per level)}

\begin{center}
\begin{tabular}{|c|c|c|c|c|c|c|c|}
\hline
\textbf{Concurrency} & \textbf{Requests} & \textbf{rps} & \textbf{p50 ms} & \textbf{p95 ms} & \textbf{p99 ms} & \textbf{Errors} & $\mathbf{X\cdot p50}$ \\ \hline
1  & 1700 & 425.0 & 2.0  & 3.8  & 5.2   & 0 & 0.8 \\
2  & 2478 & \textbf{619.5} & 2.9 & 5.0 & 6.3 & 0 & 1.8 \\
4  & 2062 & 515.5 & 7.5  & 10.9 & 13.9  & 0 & 3.9 \\
8  & 2086 & 521.5 & 15.0 & 20.8 & 25.7  & 0 & 7.8 \\
16 & 2091 & 522.8 & 29.4 & 39.6 & 85.3  & 0 & 15.4 \\
24 (extra) & 1944 & 486.0 & 46.4 & \textbf{78.1} & 153.3 & 0 & 22.6 \\
32 (extra) & 2036 & 509.0 & 60.8 & \textbf{96.5} & 129.3 & 0 & 30.9 \\ \hline
\end{tabular}
\end{center}

\begin{itemize}
  \item \textbf{Saturation.} Throughput stops growing after concurrency 2 and stays near 520\,req/s. Little's Law,
  $L = X\cdot W$ (p50 as $W$), returns the concurrency set (7.8 $\approx$ 8, 15.4 $\approx$ 16): once saturated, extra clients only queue,
  $W \approx N/X_{\max}$, and p50 doubles each time $N$ doubles. \code{top} at $c=8$ shows one uvicorn process at $\sim$92\% CPU
  (one core, the GIL) while the machine is 88\% idle. Service time $1/X_{\max} \approx 1.6$--$1.9$\,ms.
  \textbf{Saturation point: concurrency $\approx$ 2, at about 520--620 req/s.}
  \item \textbf{Capacity under p95 $\le$ 80\,ms.} With $X_{\max}\approx 520$, $W=N/520$, and p95 $\approx 1.35$--$1.7\times$ p50,
  so p95 reaches 80\,ms at p50 $\approx$ 47--59\,ms, i.e.\ $N \approx 24$--$31$. Measured: $c=24 \to$ 78.1\,ms (pass), $c=32 \to$ 96.5\,ms (fail).
  \textbf{Capacity $\approx$ 500 req/s with up to $\sim$24 requests in flight}; plan on $\sim$400\,req/s (80\%) since p99 already exceeds 80\,ms at $c=16$.
  \item \textbf{Batch 1 vs 100 ($c=1$).} Batch 1: 447\,req/s $\Rightarrow$ \textbf{2.24\,ms/row}. Batch 100: 313\,req/s $\Rightarrow$ 3.19\,ms/request
  $\Rightarrow$ \textbf{0.032\,ms/row} ($\sim$70$\times$ cheaper). Most of the cost is fixed (HTTP, JSON, pydantic, DataFrame, response); an extra row
  adds only $(3.19-2.24)/99 \approx 0.01$\,ms because Ridge is a vectorised matrix product. The price is latency: the first row waits for the last.
  \item \textbf{Why p95/p99, not the mean.} The mean hides the tail users actually see: at $c=16$, p50 is 29\,ms but p99 is 85\,ms.
  An SLO is a per-request promise, and the tail is where queueing shows first.
\end{itemize}

\textbf{Change: 2 uvicorn workers.} The CPU-bound work runs under one GIL per process, so a second process should double capacity.

\begin{itemize}
  \item At the same concurrency: \textbf{1.96$\times$ ($c=8$) and 2.02$\times$ ($c=16$)}; peak vs.\ peak $1072.8/619.5 =$ \textbf{1.73$\times$}.
  At $c=1$--$2$ nothing changes, since one request at a time cannot use two processes.
  \item p95 at $c=32$ drops from 96.5 to 40.9\,ms, so capacity under the SLO is now $\approx$ 1{,}050\,req/s. \code{top} shows both workers
  and the load generator at $\sim$100\% CPU, so the generator is now also a limit.
\end{itemize}

%================================================
\section{Part B: MLflow Registry, Gate and Shadow}

\subsection{B1. Registration (\code{partB/eta\_mlflow.py register})}
\code{EtaPyfunc} wraps the bundle and returns \code{eta\_minutes} (1 decimal) and \code{model\_version}. Each version is logged with
\code{infer\_signature} (7 float64 inputs, 2 outputs), a 2-row \code{input\_example}, pinned \code{pip\_requirements}
(scikit-learn 1.7.2, pandas 2.3.3, numpy 2.2.6, joblib 1.6.0, mlflow 3.18.0), params \code{algorithm}, \code{bundle\_version} and metric
\code{fresh\_mae}. Tracking store: \code{sqlite:///mlflow.db}.

\begin{center}
\begin{tabular}{|c|c|c|c|}
\hline
\textbf{Version} & \textbf{Bundle} & \textbf{Algorithm} & \textbf{Alias} \\ \hline
1 & v1 & Ridge (scaler + Ridge pipeline) & \textbf{prod} \\
2 & v2 & GradientBoostingRegressor & \textbf{candidate} \\ \hline
\end{tabular}
\end{center}

\subsection{B2. Promotion Gate}
\code{promotion\_gate()} evaluates \code{@prod} and \code{@candidate} on \code{make\_data(2000, seed=99)} and times 200
single-row predictions. \code{test\_gate.py}: \textbf{13 passed} (on, just inside and just outside every threshold).

\begin{center}
\begin{tabular}{|c|c|c|c|c|}
\hline
\textbf{Alias} & \textbf{MAE} & \textbf{MAE rush} & \textbf{MAE rain} & \textbf{p95 ms (in-process)} \\ \hline
prod (v1) & 11.955 & 19.680 & 11.623 & 0.90 \\
candidate (v2) & 6.105 & 8.580 & 6.285 & 1.05 \\ \hline
\end{tabular}
\end{center}
Default thresholds: relative gain $0.489 \ge 0.03$, all 4 criteria true $\Rightarrow$ \textbf{promote = true}.
With \code{min\_rel\_gain=0.60}: \code{mae\_gain} false, the rest true $\Rightarrow$ \textbf{promote = false}.

\subsection{B3. Shadow Comparison}
200 rows of \code{eta\_fresh.csv}, one row per request, the same request sent to both servers.
\begin{center}
\begin{tabular}{|l|c|}
\hline
\textbf{Metric} & \textbf{Value} \\ \hline
Mean $|\text{prod}-\text{candidate}|$ & \textbf{9.48 min} (max 51.7) \\
Requests with $|\text{diff}| > 10$ min & \textbf{35.0\%} \\
p95 latency, prod :5101 / candidate :5102 & \textbf{6.34 / 6.91 ms} (p50 4.49 / 4.68) \\ \hline
\end{tabular}
\end{center}
\begin{itemize}
  \item \textbf{Decision.} One ETA in three changes by more than 10 minutes. That is expected: v2 halves the fresh error, mostly in rush hour
  (19.7 $\to$ 8.6\,min), where Ridge cannot fit the rush $\times$ traffic interaction. Latency is almost the same and far below the SLO, and the
  gate passes on every criterion, so \textbf{promote v2} -- but through the sticky canary (10\% $\to$ 50\% $\to$ 100\% of customers) with the
  Part D abort criteria, checking MAE against real delivery times as they arrive.
  \item \textbf{Moving an alias does not change a running server.} \code{mlflow models serve -m models:/eta-model@prod} resolves the alias once at
  start-up, downloads that version and keeps it in memory; it never asks the registry again. Verified: with \code{prod} moved to version 2, the running
  :5101 server still answered \code{model\_version: v1}.
  \item \textbf{Picking up the new version without dropping requests.} (1) \emph{Rolling / blue-green restart}: start new processes (pods) that resolve the
  alias again, wait for readiness and the smoke test, switch the load balancer or Kubernetes Service, then drain the old ones
  (\code{RollingUpdate maxUnavailable: 0}, Lecture 10 Minikube). (2) \emph{Hot swap}: a background thread polls the alias (or gets a webhook), loads and
  warms the new version beside the old one, then swaps the model reference atomically; in-flight requests finish on the old model.
\end{itemize}

%================================================
\section{Part C: Ray Serve Sticky Canary}

C1 and C2: \code{partC/eta\_ray.py} (sha256 bucket of \code{customer\_id}; router 0.1 + replicas 0.25 CPU each, at most
1.35 of 2 CPUs) and \code{partC/verify\_ray.py}. All 5 checks pass (\code{partC/verify\_output.txt}): (a) 208/1000 on v2;
(b) and (c) 200/200 same arm on a second call and after a redeploy; (d) weight $0.2\to0.5$: 0 customers moved back to prod;
(e) \code{X-Beta: 1}: 50/50 on v2.

\subsection{C3. Short Answers}
\begin{enumerate}
  \item \textbf{Why not \code{hash()}?} For \code{str}, Python salts the hash per process (\code{PYTHONHASHSEED}), so every router replica and restart
  buckets customers differently -- breaking checks (b)--(d). It is also not stable across Python versions or platforms. sha256 is deterministic
  everywhere and spreads ids evenly.
  \item \textbf{Why stickiness matters.} \emph{Product:} a customer refreshing the app must not see the ETA jump between models (41 $\to$ 55 $\to$ 41 min);
  flip-flopping ETAs lose trust. \emph{Analysis:} the customer is the unit of randomisation; if one customer saw both arms, outcomes (late deliveries,
  complaints) could not be attributed to either model, and observations would not be independent. Stickiness lets delayed ground truth join back to one arm.
  \item \textbf{Autoscaling parameters.} A request takes $\sim$2--5\,ms, so \code{target\_ongoing\_requests=2} scales up as soon as a queue forms without a
  replica per request. Scale up fast (3\,s) because canary traffic is small and bursty; scale down slowly (30\,s) to avoid flapping. A 10\,s look-back with
  2\,s metrics smooths single spikes. \code{max\_replicas=3} comes from the CPU budget, and \code{min\_replicas=1} keeps a warm replica so the first canary
  customer avoids a cold start.
\end{enumerate}

%================================================
\section{Part D: Release Gate}
\code{partD/smoke\_eta.py}: 8 checks, including 3 model-sanity checks and invalid input $\to$ 422. Passes against Part A (v1, v2)
and the Ray app; exits 1 for a wrong expected version and for \code{--max-p95-ms 0.5}.

\textbf{Abort criteria for a 10\% canary of v2} (v2 arm vs.\ v1 over the same window; any one $\Rightarrow$ set weight to 0):
\begin{center}
\small
\begin{tabular}{|c|>{\raggedright}p{5.6cm}|>{\raggedright}p{4.2cm}|>{\raggedright\arraybackslash}p{4.3cm}|}
\hline
\textbf{\#} & \textbf{Metric} & \textbf{Abort if} & \textbf{Window} \\ \hline
1 & 5xx ratio of the v2 arm (\code{eta\_requests\_total\{status=\~{}"5.."\}} / all) & $> 1\%$ \textbf{or} $> 2\times$ v1's ratio & 5 min rolling, $\ge$ 200 v2 requests \\ \hline
2 & v2 p95 latency (\code{histogram\_quantile(0.95, \dots)}) & $> 80$ ms \textbf{or} $> 1.5\times$ v1 p95 & 3 consecutive 1-min windows \\ \hline
3 & MAE vs.\ real delivery times, v2 vs.\ v1 (overall, rush, rain) & v2 overall MAE $>$ v1's, or any slice $>$ v1 + 5\% & rolling 1 h, $\ge$ 300 completed v2 deliveries \\ \hline
4 & Late-delivery rate (delivered $> 15$ min after the ETA shown) & $>$ v1's rate + 2 pp & rolling 1 h, $\ge$ 300 completed v2 deliveries \\ \hline
\end{tabular}
\end{center}

%================================================
\section{Not Tested}
\begin{itemize}
  \item Autoscaling of the candidate under real load: only the config and the idle state (1 replica) were checked; no load test scaled it to 3.
  \item The abort criteria are a design only: no Prometheus server scraped \code{/metrics} (the text format was checked by tests and curl), and no alerting was set up.
  \item Client and server shared one laptop; with 2 workers the load generator was CPU-bound, so the 2-worker numbers are a lower bound.
  \item \code{/metrics} with \code{--workers 2} is per process (multiprocess mode not set up).
  \item \code{smoke\_eta.py} was not run against the MLflow servers (different API: \code{/invocations}, 400 with no range checks).
  \item The shadow test uses 200 offline rows, not mirrored live traffic; the gate's p95 is in-process \code{pyfunc.predict}, not over HTTP.
\end{itemize}

\section{Sources and Reuse}
The starter \code{eta\_data.py} and \code{load\_test.py} are used unchanged. Ideas and small fragments reused, with changes:
Lecture 10 \emph{REST\_API\_Serving\_Exercise} (\code{rest\_app.py}: bundle loaded once, pydantic schema, timing middleware,
thread-pool \code{def} endpoint, the \code{--workers} experiment); \emph{MLflow\_Serving\_Exercise} (\code{mlflow\_register.py}: pyfunc
wrapper, float64 signature, \code{log\_model} + alias; the gate on fresh data and the shadow loop from its Tasks A/B); \emph{Ray\_Serve\_Exercise}
(\code{ray\_serve\_app.py}: model deployment + FastAPI Router + \code{build\_app}, fractional CPUs, \code{extra\_work\_ms}, autoscaling keys);
Lecture 10 \emph{Minikube\_Serving} slides (smoke-test checks and \code{smoke || rollback}, rolling update for B3); Lecture 9 and my
Assignment 4 (Prometheus counter and histogram).

\end{document}
